\documentclass[12pt]{article}
\usepackage{styles/plantilla-profesor}

\usepackage{graphicx}
\usepackage[spanish]{babel}

\setlength{\parskip}{8pt}   % Espacio entre párrafos normales

\usepackage{titlesec}
\titlespacing*{\section}{0pt}{10pt}{5pt}    % antes, después
\titlespacing*{\subsection}{0pt}{8pt}{4pt}
\titlespacing*{\subsubsection}{0pt}{6pt}{3pt}

\begin{document}

\begin{titlepage}
    \centering
    {\LARGE Universidad Nacional Mayor de San Marcos \par}
    {\large Universidad del Perú. Decana de América\par}
    \vspace{0.5cm}
    {\Large Facultad de Ingeniería de Sistemas e Informática\par}
    \vspace{0.5cm}

    \includegraphics[width=8cm]{imagenes/logo-unmsm.png}\par
    \vspace{0.5cm}

    {\LARGE\bfseries Asignación Óptima de Recursos y Tareas en Metodologías Ágiles\par}
    \vspace{0.5cm}
    {\large\bfseries Aplicación de Algoritmos Genéticos en Entornos SCRUM\par}
    \vspace{1cm}

    {\Large\bfseries Integrantes:\par}
    \vspace{0.5cm}
    Salazar Garcia, Diego\\[4pt]
    Arias Chumpitaz, Giovanni\\[4pt]
    Arroyo Vasquez, Luis\\[4pt]
    Lavaud Guevara, Jean\\[4pt]
    \vspace{1cm}

    {\large \textbf{Docente:} Claudio Arango\par}
    \vfill

    {\large Lima, Perú\par}
    {\large 2025\par}
\end{titlepage}


\section*{Resumen}
Este informe presenta una solución basada en algoritmos genéticos para abordar el problema de congestión vehicular en Lima. Se describe el contexto del tráfico urbano y la aplicación de técnicas de optimización bioinspiradas.


\section{Problema}
En entornos ágiles como SCRUM, una de las dificultades recurrentes es la asignación eficiente de tareas entre los miembros del equipo, considerando factores como carga de trabajo, tiempo estimado, prioridades y dependencias.
\vspace{0.5cm}

\begin{tcolorbox}[colback=gray!10, colframe=black!30, title={Sugerencia para esta sección}]
    Describe el problema real: cómo se distribuyen mal las tareas, cómo puede haber sobrecarga en algunos programadores y cómo las dependencias entre tareas generan cuellos de botella.
\end{tcolorbox}

\section*{Objetivos}
\begin{itemize}
  \item Desarrollar un algoritmo genético para optimizar rutas.
  \item Simular su comportamiento bajo distintas condiciones urbanas.
  \item Validar los resultados con datos reales de tráfico.
\end{itemize}

\section{Alcance del Proyecto}
Este proyecto se enfocará en la asignación de tareas para un equipo ágil ficticio conformado por desarrolladores, analistas y testers. No se abordarán aspectos de gestión externa o requerimientos no funcionales.
\vspace{0.5cm}

\begin{tcolorbox}[colback=gray!10, colframe=black!30, title={Sugerencia para esta sección}]
    Define claramente qué incluye tu solución (asignación interna de tareas, esfuerzo estimado, simulación SCRUM) y qué no (por ejemplo, herramientas reales, despliegue).
\end{tcolorbox}


\section{Marco Teórico}

Según \textcite{scrum2020}, la asignación de tareas dentro de metodologías ágiles como SCRUM requiere un entendimiento profundo de conceptos relacionados con planificación de proyectos, inteligencia artificial y optimización. Los algoritmos genéticos, inspirados en los procesos evolutivos naturales, ofrecen una alternativa poderosa para resolver problemas complejos con múltiples variables y restricciones \parencite{goldberg}.

% … o, si quieres el apellido dentro del texto:
% Según \textcite{goldberg}, los algoritmos genéticos ofrecen …

Dentro del marco de este trabajo, se exploran teorías relacionadas con el problema del viajante (TSP), la gestión ágil de proyectos, técnicas de optimización y fundamentos de los algoritmos genéticos aplicados a software.

\vspace{0.5cm}

\begin{tcolorbox}[colback=gray!10,colframe=black!30,title={Sugerencias para esta sección}]
    \begin{itemize}
        \item Introduce brevemente conceptos como SCRUM, DevOps, optimización, TSP, etc.
        \item Cita autores relevantes con formato APA 7 si es necesario.
        \item Relaciona estos conceptos con el problema que estás abordando.
        \item Puedes usar gráficos teóricos o tablas si fuera necesario.
    \end{itemize}
\end{tcolorbox}

\section*{3. Metodología}
Se sigue un enfoque incremental: definición de objetivos, modelado del TSP, diseño del algoritmo genético y validación. Se incluye flujograma y arquitectura de la solución.

\section{Propuesta de Solución}

La propuesta presentada consiste en un sistema interactivo orientado a entornos ágiles, que permite optimizar la asignación de tareas dentro de un equipo de desarrollo mediante el uso de un algoritmo genético multiobjetivo. La solución busca balancear la carga de trabajo, minimizar el tiempo total de ejecución y considerar factores clave como dependencias, habilidades requeridas y costos asociados.

El sistema se estructura en torno a una arquitectura modular que separa claramente la interfaz de usuario, la lógica de negocio y el motor de optimización. Esta separación permite una interacción clara y eficiente con el usuario, al mismo tiempo que encapsula la complejidad de la heurística evolutiva implementada.

En las siguientes subsecciones se detalla el modelo general del sistema, el diseño del algoritmo genético propuesto, y la forma en que los componentes se integran para generar una solución viable y extensible al problema de asignación de tareas.

\subsection{Modelo General del Sistema}

Como se representó anteriormente en la Figura~\ref{fig:arquitectura}, el sistema está compuesto por tres componentes principales: \textbf{Interfaz Gráfica}, \textbf{API} y \textbf{Núcleo de Optimización}.

\begin{itemize}
    \item \textbf{Interfaz Gráfica}: Permite al usuario visualizar los resultados y configurar parámetros clave del algoritmo genético, como el tamaño de la población, número de generaciones, tasas de cruce y mutación, tamaño del torneo y los pesos de la función objetivo.

    \item \textbf{API}: Recibe las configuraciones del usuario y los datos estáticos de tareas y desarrolladores. Actúa como intermediario entre la interfaz gráfica y el núcleo del algoritmo, enviando y recibiendo información a través de endpoints definidos en Flask.

    \item \textbf{Núcleo de Optimización}: Implementado en Python, este módulo contiene la lógica completa del algoritmo genético. Se encarga de generar soluciones viables y optimizadas para la asignación de tareas, respetando restricciones de dependencias, habilidades, tiempos y costos.
\end{itemize}

Esta aproximación cliente-servidor con separación clara entre capas permite al usuario interactuar de forma sencilla con la aplicación, facilitando la configuración de parámetros y la visualización de resultados, todo mientras se mantiene encapsulada la complejidad interna del algoritmo genético multiobjetivo para la optimización de la asignación de tareas.

\subsection{Algoritmo Genético Propuesto}

El núcleo de optimización sigue un enfoque evolutivo, donde una población inicial de soluciones (cromosomas) se somete a iteraciones de evaluación, selección, cruce y mutación para mejorar la asignación de tareas.

\begin{itemize}
    \item Se parte de una población aleatoria de asignaciones.
    \item Cada generación aplica operadores genéticos y se evalúan las soluciones.
    \item La mejor solución se toma como sprint generado.
    \item El proceso se repite hasta vaciar el backlog.
\end{itemize}

\noindent\textbf{Imagen:} \textit{(opcional)} flujograma simple del algoritmo. Ej: pseudodiagrama basado en 'while tareas\_pendientes:'.

\subsection{Representación de Soluciones}

Cada cromosoma representa una posible asignación de tareas a desarrolladores. Esta representación codifica:

\begin{itemize}
    \item Qué tareas se asignan a qué programadores.
    \item El orden en que deben completarse (respetando dependencias).
    \item El ajuste según las habilidades de cada desarrollador.
\end{itemize}

\noindent\textbf{Imagen:} \textit{(opcional)} ejemplo visual de un cromosoma o tabla de asignaciones (opcional si ya lo describes en texto).

\subsection{Evaluación de Soluciones (Función Objetivo)}

La función objetivo pondera múltiples criterios para determinar la calidad de una solución:

\begin{itemize}
    \item \textbf{Makespan (50\%)}: duración del sprint.
    \item \textbf{Balance de carga (25\%)}: varianza en horas por developer.
    \item \textbf{Coincidencia de habilidades (20\%)}: afinidad entre habilidades requeridas y ofrecidas.
    \item \textbf{Costo (5\%)}: según tarifas por hora.
\end{itemize}

\noindent\textbf{Imagen:} \textit{(opcional)} fórmula matemática o pseudocódigo de evaluación.

\subsection{Operadores Genéticos Personalizados}

Los operadores han sido ajustados para respetar restricciones del problema:

\begin{itemize}
    \item \textbf{Selección:} por torneo.
    \item \textbf{Cruce:} combinación parcial de asignaciones entre padres.
    \item \textbf{Mutación:} reasignación de tareas con baja probabilidad.
    \item \textbf{Corrección:} validación para evitar duplicados y tareas inválidas.
\end{itemize}

\noindent\textbf{Imagen:} no requerida, puede ser solo texto.

\subsection{Interfaz y Configuración del Usuario}

El sistema permite que el usuario ajuste los parámetros del algoritmo genético desde una interfaz web amigable:

\begin{itemize}
    \item Población inicial
    \item Número de generaciones
    \item Tasa de mutación y cruce
    \item Tamaño del torneo
    \item Pesos para cada criterio de evaluación
\end{itemize}

\noindent\textbf{Imagen:} captura del frontend (panel de configuración).

\subsection{Ejemplo de Ejecución}

Se realizó una ejecución completa con el conjunto estático de tareas y desarrolladores. Se configuraron los parámetros del algoritmo y se generaron varios sprints:

\begin{itemize}
    \item Se obtuvo un conjunto de sprints con carga balanceada.
    \item Las tareas respetaron todas sus dependencias.
    \item El makespan fue razonable para el conjunto dado.
\end{itemize}

\noindent\textbf{Imagen:} \textit{opcional} captura de los resultados visualizados (por ejemplo, tabla con desarrolladores y tareas).


% La propuesta consiste en un sistema capaz de asignar automáticamente tareas a desarrolladores en base a un modelo evolutivo. El sistema considera el esfuerzo estimado por tarea, las dependencias entre ellas y la disponibilidad de los miembros del equipo.

% El algoritmo genético es ajustado para minimizar el desequilibrio de carga, respetar las restricciones impuestas por los sprints y optimizar el flujo de trabajo. La solución incluye un módulo de visualización que muestra la distribución sugerida.

% \vspace{0.5cm}

% \begin{tcolorbox}[colback=gray!10, colframe=black!30, title={Sugerencias para esta sección}]
%     \begin{itemize}
%         \item Explica brevemente qué hace tu solución y cómo se integra en un entorno ágil.
%         \item Muestra capturas o diagramas si ya has hecho avances de código.
%         \item Puedes separar en módulos: asignación, evaluación, visualización.
%         \item Justifica por qué tu propuesta mejora lo que ya existe.
%     \end{itemize}
% \end{tcolorbox}
\section{Conclusiones}

Este trabajo ha demostrado que es posible construir un sistema automatizado capaz de resolver la asignación de tareas en entornos ágiles SCRUM mediante un algoritmo genético multiobjetivo, integrando simultáneamente múltiples restricciones prácticas como habilidades requeridas, dependencias entre tareas, capacidades individuales y costos. La solución no solo logra generar asignaciones eficientes, sino que también respeta principios clave de metodologías ágiles como la iteración progresiva y la transparencia del proceso.

La implementación completa, compuesta por una interfaz gráfica en React, una API desarrollada en Flask y un motor de optimización en Python, valida que un enfoque modular y desacoplado es viable incluso en proyectos académicos, sentando las bases para futuras integraciones más complejas. Este diseño permitió realizar validaciones constantes a lo largo del desarrollo, ajustando y refinando componentes del algoritmo hasta alcanzar una versión estable y funcional.

Uno de los principales aprendizajes fue comprobar que, con una buena representación del problema y operadores genéticos bien diseñados, es posible obtener soluciones de alta calidad incluso en un dominio con restricciones duras y objetivos conflictivos. También se confirmó que la exposición clara de los parámetros de configuración al usuario (como tasas de cruce, mutación o pesos de la función objetivo) agrega valor al sistema, haciéndolo más flexible y comprensible.

Si bien existen limitaciones como el uso de datos estáticos para tareas y desarrolladores, estas no comprometen el valor de la propuesta como prueba de concepto sólida. El sistema desarrollado no solo cumple su objetivo académico, sino que representa un punto de partida válido para aplicaciones reales en entornos empresariales o educativos donde se requiere automatizar la gestión de tareas de forma inteligente.

% \section{Conclusiones}

% El enfoque basado en algoritmos genéticos ha demostrado ser una herramienta eficiente para la asignación de tareas en entornos SCRUM. Se ha logrado un balance adecuado en la distribución del trabajo, respetando restricciones y optimizando el tiempo total del proyecto.

% Este trabajo sienta las bases para aplicaciones reales en la industria del software, y abre la posibilidad de integrar técnicas similares en procesos más amplios de gestión ágil.

% \vspace{0.5cm}

% \begin{tcolorbox}[colback=gray!10, colframe=black!30, title={Sugerencias para esta sección}]
%     \begin{itemize}
%         \item Resume los aprendizajes clave del trabajo realizado.
%         \item Redacta al menos 3 conclusiones claras y bien separadas si es posible.
%         \item Puedes mencionar limitaciones y oportunidades de mejora.
%         \item Evita repetir el resumen o los objetivos textualmente.
%     \end{itemize}
% \end{tcolorbox}
\section{Recomendaciones}

Una de las principales recomendaciones es extender la solución para permitir la carga dinámica de datos por parte del usuario. Actualmente, las tareas y desarrolladores están predefinidos en el sistema, lo cual limita su flexibilidad. Incorporar formularios o integración con bases de datos permitiría simular escenarios reales y adaptar el sistema a contextos específicos de empresas o instituciones educativas. Esta mejora también facilitaría la validación del algoritmo con conjuntos de datos reales y facilitaría su adopción práctica.

Adicionalmente, se sugiere mejorar las capacidades de visualización del sistema. Aunque la interfaz actual permite observar resultados básicos, sería conveniente integrar herramientas interactivas como gráficos de Gantt, dashboards con métricas de rendimiento o incluso simulaciones visuales del avance del sprint. Estas mejoras no solo enriquecerían la experiencia del usuario, sino que también permitirían tomar decisiones más informadas a partir de los resultados generados por el algoritmo.

Finalmente, desde una perspectiva investigativa, se recomienda experimentar con variantes del algoritmo genético e incluso explorar enfoques híbridos, como algoritmos meméticos o combinaciones con reglas heurísticas extraídas de prácticas ágiles reales. También sería interesante aplicar el sistema en casos de estudio reales, con equipos de desarrollo reales, para evaluar no solo métricas técnicas como el makespan o la varianza de carga, sino también el impacto en la motivación, percepción de equidad y productividad del equipo.

% \section{Recomendaciones}

% A partir del análisis y la implementación realizada, se proponen las siguientes recomendaciones para futuras investigaciones o mejoras en sistemas SCRUM asistidos por algoritmos genéticos:

% \begin{itemize}
%     \item Integrar herramientas de visualización para facilitar la interpretación de las asignaciones realizadas por el algoritmo.
%     \item Explorar técnicas híbridas que combinen algoritmos genéticos con heurísticas basadas en reglas de negocio.
%     \item Validar el modelo con datos reales de proyectos ágiles en diferentes industrias.
%     \item Evaluar el impacto del algoritmo en la satisfacción y productividad del equipo.
% \end{itemize}

% \vspace{0.5cm}

% \begin{tcolorbox}[colback=gray!10, colframe=black!30, title={Sugerencia para esta sección}]
%     Formula recomendaciones específicas, factibles y que se deriven directamente de los resultados de tu propuesta. Considera aspectos como mejoras técnicas, escalabilidad o nuevas direcciones de investigación.
% \end{tcolorbox}

\newpage
\section{Bibliografía}

\begin{thebibliography}{99}

  \bibitem{scrum-guide}
  Schwaber, K. \& Sutherland, J. (2020). \textit{The Scrum Guide}. Scrum.org. Disponible en: \url{https://scrumguides.org}

  \bibitem{goldberg}
  Goldberg, D. E. (1989). \textit{Genetic Algorithms in Search, Optimization and Machine Learning}. Addison-Wesley.

  \bibitem{devops}
  Humble, J., \& Farley, D. (2010). \textit{Continuous Delivery: Reliable Software Releases through Build, Test, and Deployment Automation}. Addison-Wesley.

  \bibitem{apa7}
  American Psychological Association. (2020). \textit{Publication Manual of the American Psychological Association} (7th ed.).

\end{thebibliography}

\vspace{0.5cm}

\begin{tcolorbox}[colback=gray!10, colframe=black!30, title={Sugerencia para esta sección}]
  Usa el estilo APA 7 para citar fuentes académicas. Puedes utilizar herramientas como Zotero, Mendeley o Google Scholar para exportar referencias en BibTeX y usarlas con el paquete \texttt{biblatex} si deseas automatizar esta sección.
\end{tcolorbox}

\newpage

\section*{Autores}
% TODO: Investigar mas acerca de archivos .bib, y librerias de gestion de referencias para apa7

\begin{itemize}
  % \includegraphics[width=1.5cm]{imagenes/diego.jpg}
  \item \textbf{Salazar Garcia, Diego}: Estudiante de ingeniería con interés en optimización y sistemas inteligentes.

        % \includegraphics[width=1.5cm]{imagenes/giovanni.jpg}
  \item \textbf{Arias Chumpitaz, Giovanni}: Apasionado por el análisis de datos y metodologías ágiles.

        % \includegraphics[width=1.5cm]{imagenes/luis.jpg}
  \item \textbf{Arroyo Vasquez, Luis}: Interesado en automatización de procesos y calidad de software.

        % \includegraphics[width=1.5cm]{imagenes/jean.jpg}
  \item \textbf{Lavaud Guevara, Jean}: Enfocado en algoritmos evolutivos y arquitectura ágil de proyectos.
\end{itemize}

\vspace{0.5cm}

\begin{tcolorbox}[colback=gray!10, colframe=black!30, title={Sugerencia para esta sección}]
  Incluye una breve descripción de cada autor, como sus intereses académicos o técnicos. Puedes usar fotos pequeñas con \texttt{\\includegraphics} si lo deseas más adelante.
\end{tcolorbox}


\end{document}


% \documentclass{article}
% \usepackage{styles/plantilla-profesor} % <-- Aquí se importa la plantilla

% \begin{document}

% \section*{Prueba de la plantilla del profesor}

% Este es un pequeño ejemplo para comprobar que todo está funcionando correctamente, incluso con la plantilla en el archivo styles. Aquí mostramos una fórmula matemática:

% \[
%   f(x) = x^2 + 3x + 4
% \]

% Y también un ejemplo de código con estilo personalizado:

% \begin{lstlisting}[style=matlabstyle, language=Python, caption={Algoritmo simple}]
% def cuadrado(x):
%     return x * x

% print(cuadrado(5))
% \end{lstlisting}

% \begin{tcolorbox}[contexto]
%   Este recuadro representa un contexto teórico que podrías usar en tu informe.
% \end{tcolorbox}

% \begin{tcolorbox}[ejercicio]
%   Desarrolla un experimento para verificar el rendimiento de tu algoritmo genético con distintos tamaños de población.
% \end{tcolorbox}

% \begin{tcolorbox}[conclusion]
%   La plantilla se ha compilado correctamente. Puedes comenzar a escribir tu informe a partir de esta base.
% \end{tcolorbox}

% \end{document}
